\documentclass[a4paper,11pt]{article}

\usepackage{../../../utils/template}

\newgeometry{ left=1.0cm, right=1.0cm, top=0.8cm, bottom=0.8cm }

\pagestyle{empty}

% ===== 题目集版式宏 =====
% 章标题：独立成页大标题
\newcommand{\chaptermark}[1]{%
    \newpage
    \begin{center}
        {\Large\textbf{#1}}
    \end{center}
    \vspace{1em}}
% 节标题
\newcommand{\booksec}[1]{\vspace{0.55em}\noindent\textbf{#1}\vspace{0.2em}\par}
% 习题块标题（如「习题 1-1」）
\newcommand{\exgroup}[1]{\vspace{0.6em}\noindent\textbf{#1}\par\vspace{0.2em}}
% 例题：\examp{例 x.x}{题干}
\newcommand{\examp}[2]{\vspace{0.65em}\noindent\textbf{#1}\quad #2\par\vspace{0.3em}}
% 习题列表
\newenvironment{probs}
    {\begin{enumerate}[leftmargin=2.2em, itemsep=0.55em]}
    {\end{enumerate}}

% ===== 选项排版（均靠左）=====
% opts1：每个选项独占一行（长选项）
\newenvironment{opts1}
    {\par\vspace{0.2em}\noindent\renewcommand{\arraystretch}{1.25}\begin{tabular}{@{}p{0.96\linewidth}@{}}}
    {\end{tabular}\par\vspace{0.2em}}
% opts2：两个选项一行，各占 50%
\newenvironment{opts2}
    {\par\vspace{0.2em}\noindent\renewcommand{\arraystretch}{1.25}\begin{tabular}{@{}p{0.475\linewidth}p{0.475\linewidth}@{}}}
    {\end{tabular}\par\vspace{0.2em}}
% opts4：四个选项一行，各占 25%
\newenvironment{opts4}
    {\par\vspace{0.2em}\noindent\renewcommand{\arraystretch}{1.25}\begin{tabular}{@{}p{0.235\linewidth}p{0.235\linewidth}p{0.235\linewidth}p{0.235\linewidth}@{}}}
    {\end{tabular}\par\vspace{0.2em}}

\begin{document}

% 注意：本文件由 main_random.tex 生成（顺序与其完全一致，main_random 由 main.tex seed=20261127 打乱）。
% 例x.x 题号已替换为序号 1..205；每题上方注释保留「序号 = 例x.x」对照，PDF 不显示。
% 题目改动请改 main.tex，再重新生成 main_random.tex 和本文件，勿直接改。

% 1 = 例13.14
\vbox{%
\examp{1.}{设 $\dfrac{(x+ay)\,\mathrm{d}x+y\,\mathrm{d}y}{(x+y)^2}$ 是某个二元函数的全微分，求 $a$ 的值.}
}

% 2 = 例11.9
\vbox{%
\examp{2.}{设数列 $\{a_{n}\}$ 的通项 $a_{n}=\displaystyle\int_{0}^{+\infty}\dfrac{\mathrm{d}x}{(1+x^{2})^{n}}$，$n=2,3,\cdots$，计算 $\lim\limits_{n\to\infty}\left(\dfrac{a_{n+1}}{a_{n}}\right)^{\ln(1+\mathrm{e}^{2n})}$.}
}

% 3 = 例4.9
\vbox{%
\examp{3.}{已知函数 $f(x)=\mathrm{e}^{x}+2x+1$，设 $g(y)$ 与 $f(x)$ 互为反函数，则 $g''(2)=$（\quad）.}
\begin{opts4}
(A) $\frac{1}{3}$ & (B) $-3$ & (C) $-\frac{1}{27}$ & (D) $-\frac{\mathrm{e}^{2}}{(\mathrm{e}^{2}+2)^{3}}$
\end{opts4}
}

% 4 = 例13.8
\vbox{%
\examp{4.}{已知函数 $f(x,y)=\begin{cases}\dfrac{xy}{x^2+y^2}, & (x,y)\neq(0,0),\\ 0, & (x,y)=(0,0),\end{cases}$ 则 $f(x,y)$ 在点 $(0,0)$ 处（\quad）.}
\begin{opts2}
(A) 连续，偏导数存在 & (B) 连续，偏导数不存在 \\
(C) 不连续，偏导数存在 & (D) 不连续，偏导数不存在
\end{opts2}
}

% 5 = 例11.1
\vbox{%
\examp{5.}{求极限 $\lim\limits_{h\to 0^{+}}\displaystyle\int_{-1}^{1}\dfrac{h}{h^{2}+x^{2}}f(x)\,\mathrm{d}x$，其中函数 $f(x)$ 在 $[-1,1]$ 上连续.}
}

% 6 = 例8.2
\vbox{%
\examp{6.}{设 $a\in(0,1)$，则 $\lim\limits_{n\to\infty}\sum\limits_{i=1}^{n}\dfrac{i}{ni+a}\sin\dfrac{i}{n}=$（\quad）.}
\begin{opts4}
(A) $a-\cos a$ & (B) $a-\sin a$ & (C) $1-\cos 1$ & (D) $1-\sin 1$
\end{opts4}
}

% 7 = 例12.4
\vbox{%
\examp{7.}{如图所示，一闸门的上部是一个宽为 $2\,\mathrm{m}$、高为 $H\,\mathrm{m}$ 的矩形，下部由 $y=x^2$ 与 $y=1$ 围成. 当闸门上边缘与水面在一个平面时，其上部所受水压力与下部所受水压力之比为 $\dfrac{5}{4}$，求上部的高度 $H$.}
}

% 8 = 例9.8
\vbox{%
\examp{8.}{$\displaystyle\int\mathrm{e}^x\left(\dfrac{1-x}{1+x^2}\right)^{2}\mathrm{d}x=\underline{\hspace{3em}}$.}
}

% 9 = 例9.22
\vbox{%
\examp{9.}{$\displaystyle\int_0^1\sqrt{2x-x^2}\,\mathrm{d}x=\underline{\hspace{3em}}$.}
}

% 10 = 例11.19
\vbox{%
\examp{10.}{设 $f(x)$ 在 $[0,1]$ 上具有二阶连续导数，$f(0)=f(1)=0$，$f''(x)<0$，证明：对于任意 $x\in(0,1)$，$\displaystyle\int_0^1 \left|\dfrac{f''(x)}{f(x)}\right|\mathrm{d}x \geqslant 4$.}
}

% 11 = 例11.4
\vbox{%
\examp{11.}{$I=\displaystyle\int_{-1}^{1}\dfrac{\mathrm{d}x}{1+\mathrm{e}^{\frac{1}{x}}}=\underline{\hspace{3em}}$.}
}

% 12 = 例14.1
\vbox{%
\examp{12.}{$\lim\limits_{n\to\infty}\sum\limits_{i=1}^{n}\sum\limits_{j=1}^{n}\dfrac{i}{(n+i)(n^{2}+j^{2})}=\underline{\hspace{3em}}$.}
}

% 13 = 例9.1
\vbox{%
\examp{13.}{求 $\displaystyle\int\dfrac{x}{x^2+2x+2}\,\mathrm{d}x$.}
}

% 14 = 例15.14
\vbox{%
\examp{14.}{设曲线 $L\colon y=y(x)$（$x>\mathrm{e}$）经过点 $(\mathrm{e}^{2},0)$，$L$ 上任一点 $P(x,y)$ 到 $y$ 轴的距离等于该点处的切线在 $y$ 轴上的截距. \\
(1) 求 $y(x)$；\\
(2) 在 $L$ 上求一点，使该点处的切线与两坐标轴所围三角形的面积最小，并求此最小面积.}
}

% 15 = 例1.3
\vbox{%
\examp{15.}{设 $f(x)=\begin{cases}\lim\limits_{t\to x}\left(\dfrac{x-1}{t-1}\right)^{\frac{t}{x-t}}, & x\neq 1,\\[1ex] 0, & x=1,\end{cases}$ 则 $f(x)$ 的第二类间断点的个数为（\quad）.}
\begin{opts4}
(A) $0$ & (B) $1$ & (C) $3$ & (D) $\infty$
\end{opts4}
}

% 16 = 例2.11
\vbox{%
\examp{16.}{设 $u=g(x)=\begin{cases} x, & x \text{ 是有理数} \\ 0, & x \text{ 是无理数} \end{cases}$，$y=f(u)=\begin{cases} 1, & u\neq 0 \\ 0, & u=0 \end{cases}$，则 $\lim\limits_{x\to x_0}f[g(x)]$（\quad）.}
\begin{opts4}
(A) 等于 $0$ & (B) 等于 $1$ & (C) 等于 $x$ & (D) 不存在
\end{opts4}
}

% 17 = 例7.1
\vbox{%
\examp{17.}{质点在第一象限沿曲线 $y=x^{\frac{3}{2}}$ 远离原点，已知质点与原点的距离以 $11\,\mathrm{cm/s}$ 的变化率增加，则当质点运动到 $x=3\,\mathrm{cm}$ 时的水平速度为\underline{\hspace{3em}}.}
}

% 18 = 例1.2
\vbox{%
\examp{18.}{计算 $\lim\limits_{x\to 0}\left[\dfrac{1+\displaystyle\int_{0}^{x}\dfrac{\sin t}{t}\,\mathrm{d}t}{x}-\dfrac{1}{\ln(1+x)}\right]$.}
}

% 19 = 例5.1
\vbox{%
\examp{19.}{设 $f(x)$ 有连续的一阶导数，且 $f(0)=0$，$f'(0)=1$. 求极限 $\lim\limits_{x\to 0}\dfrac{xf(u)}{uf(x)}$，其中 $u$ 是曲线 $y=f(x)$ 在点 $(x,f(x))$ 处的切线在 $x$ 轴上的截距.}
}

% 20 = 例9.10
\vbox{%
\examp{20.}{计算 $\displaystyle\int_0^1\dfrac{1}{(x+1)(x^2-2x+2)}\,\mathrm{d}x$.}
}

% 21 = 例6.21
\vbox{%
\examp{21.}{设函数 $y=y(x)$ 由参数方程 $\begin{cases} x=\dfrac{1}{3}t^3+t,\\[1ex] y=\dfrac{1}{3}t^3-t+\dfrac{1}{3} \end{cases}$ 确定. \\
(1) 求 $y(x)$ 的极值；\\
(2) 若 $\begin{cases} u=x,\\ v=y+k, \end{cases}$ 且 $v=v(u)$ 恰有一个零点，求常数 $k$ 的取值范围.}
}

% 22 = 例6.2
\vbox{%
\examp{22.}{已知函数 $f(x)$ 和 $g(x)$ 在 $[a,b]$ 上连续，在 $(a,b)$ 内可导，且 $g'(x)\neq 0$，证明：存在 $\xi\in(a,b)$，使得 $\dfrac{f(\xi)-f(a)}{g(b)-g(\xi)}=\dfrac{f'(\xi)}{g'(\xi)}$.}
}

% 23 = 例9.11
\vbox{%
\examp{23.}{求下列不定积分：\\[0.6ex]
(1) $\displaystyle\int\dfrac{\sin^5 x}{\cos^4 x}\,\mathrm{d}x$；\\[0.6ex]
(2) $\displaystyle\int\dfrac{1}{\sin x+\sin^3 x}\,\mathrm{d}x$.}
}

% 24 = 例15.15
\vbox{%
\examp{24.}{求一条凹曲线 $y=y(x)$（$x\geqslant 1$）的表达式，已知其上任一点处的曲率 $k=\dfrac{1}{2y^{2}\cos\alpha}$，其中 $\alpha$ 为该曲线在相应点处的切线的倾角，$\cos\alpha>0$，并设曲线在点 $(3,2)$ 处的切线的倾角为 $45^{\circ}$.}
}

% 25 = 例6.10
\vbox{%
\examp{25.}{设函数 $f(x)$ 在区间 $[a,b]$ 上满足：对任意 $x$, $y\in[a,b]$，有 $|f(x)-f(y)|\leqslant M|x-y|^{\alpha}$，其中 $M>0$，$\alpha>1$ 是常数. 证明：$f(x)$ 在 $[a,b]$ 上恒为常数.}
}

% 26 = 例9.20
\vbox{%
\examp{26.}{函数 $f(x)=\begin{cases}\dfrac{1}{\sqrt{1+x^2}}, & x\leqslant 0,\\ (x+1)\cos x, & x>0\end{cases}$ 的一个原函数为（\quad）.}
\begin{opts1}
(A) $F(x)=\begin{cases}\ln(\sqrt{1+x^2}-x), & x\leqslant 0,\\ (x+1)\cos x-\sin x, & x>0\end{cases}$ \\[0.8ex]
(B) $F(x)=\begin{cases}\ln(\sqrt{1+x^2}-x)+1, & x\leqslant 0,\\ (x+1)\cos x-\sin x, & x>0\end{cases}$ \\[0.8ex]
(C) $F(x)=\begin{cases}\ln(\sqrt{1+x^2}+x), & x\leqslant 0,\\ (x+1)\sin x+\cos x, & x>0\end{cases}$ \\[0.8ex]
(D) $F(x)=\begin{cases}\ln(\sqrt{1+x^2}+x)+1, & x\leqslant 0,\\ (x+1)\sin x+\cos x, & x>0\end{cases}$
\end{opts1}
}

% 27 = 例2.2
\vbox{%
\examp{27.}{若 $x_{n+1}=\dfrac{\pi}{2}+\dfrac{1}{2}\cos x_{n}$（$n=1,2,\cdots$），$x_{1}=\pi$，证明数列 $\{x_{n}\}$ 收敛，并求 $\lim\limits_{n\to\infty}x_{n}$.}
}

% 28 = 例13.3
\vbox{%
\examp{28.}{求 $\lim\limits_{\substack{x\to\infty\\ y\to\infty}}\dfrac{x+y}{x^2-xy+y^2}$.}
}

% 29 = 例2.12
\vbox{%
\examp{29.}{已知数列 $\{x_n\}$，其中 $-\dfrac{\pi}{2}\leqslant x_n\leqslant\dfrac{\pi}{2}$，则（\quad）.}
\begin{opts1}
(A) 若 $\lim\limits_{n\to\infty}\cos\sin x_n$ 存在，则 $\lim\limits_{n\to\infty}x_n$ 存在\\
(B) 若 $\lim\limits_{n\to\infty}\sin\cos x_n$ 存在，则 $\lim\limits_{n\to\infty}x_n$ 存在\\
(C) 若 $\lim\limits_{n\to\infty}\cos\sin x_n$ 存在，则 $\lim\limits_{n\to\infty}\sin x_n$ 存在，但 $\lim\limits_{n\to\infty}x_n$ 不一定存在\\
(D) 若 $\lim\limits_{n\to\infty}\sin\cos x_n$ 存在，则 $\lim\limits_{n\to\infty}\cos x_n$ 存在，但 $\lim\limits_{n\to\infty}x_n$ 不一定存在
\end{opts1}
}

% 30 = 例13.6
\vbox{%
\examp{30.}{设 $f(x,y)=\dfrac{x^2-y^2}{x^2+y^2}$，$I_1=\lim\limits_{x\to 0}\left[\lim\limits_{y\to 0}f(x,y)\right]$，$I_2=\lim\limits_{y\to 0}\left[\lim\limits_{x\to 0}f(x,y)\right]$，$I_3=\lim\limits_{\substack{x\to 0\\ y\to 0}}f(x,y)$，则（\quad）.}
}

% 31 = 例1.1
\vbox{%
\examp{31.}{$\lim\limits_{x\to 0^{+}}\dfrac{x^{x}-1}{\ln x\cdot\ln(1-x)}=\underline{\hspace{3em}}$.}
}

% 32 = 例13.5
\vbox{%
\examp{32.}{设 $f(x,y)=\dfrac{xy}{x^2+y^2}$，$I_1=\lim\limits_{\substack{x\to 0\\ y\to 0}}f(x,y)$，$I_2=\lim\limits_{y\to 0}\left[\lim\limits_{x\to 0}f(x,y)\right]$，则（\quad）.}
\begin{opts2}
(A) $I_1$，$I_2$ 均存在 & (B) $I_1$ 存在，$I_2$ 不存在 \\
(C) $I_1$ 不存在，$I_2$ 存在 & (D) $I_1$，$I_2$ 均不存在
\end{opts2}
}

% 33 = 例13.30
\vbox{%
\examp{33.}{求 $u=\sqrt{x^2+y^2}$ 在约束条件 $5x^2+4xy+2y^2=1$ 下的最大值与最小值.}
}

% 34 = 例13.32
\vbox{%
\examp{34.}{设函数 $f(u)$ 在区间 $(0,+\infty)$ 内具有三阶连续导数，且存在可微函数 $F(x,y)$ 使得 $\mathrm{d}[F(x,y)]=\dfrac{f(xy)}{x^2 y}\,\mathrm{d}x+\dfrac{f''(xy)}{xy^2}\,\mathrm{d}y$（$xy>0$）. \\
(1) 证明 $\dfrac{f''(u)}{u}-\dfrac{f(u)}{u}=C$，$C$ 为常数；\\
(2) 设 $f(1)=1$，$f'(1)=-1$，$f''(1)=0$，求 $f(u)$ 的表达式.}
}

% 35 = 例9.7
\vbox{%
\examp{35.}{求下列不定积分：\\[0.6ex]
(1) $\displaystyle\int\mathrm{e}^x\sin x\,\mathrm{d}x$；\\[0.6ex]
(2) $\displaystyle\int\sqrt{1+x^2}\,\mathrm{d}x$.}
}

% 36 = 例9.18
\vbox{%
\examp{36.}{设函数 $f(x)=\displaystyle\int_0^1|t^2-x^2|\,\mathrm{d}t$（$x>0$），求 $f'(x)$，并求 $f(x)$ 的最小值.}
}

% 37 = 例10.1
\vbox{%
\examp{37.}{求心形线 $r=a(1-\cos\theta)$（$a>0$，$0\leqslant\theta\leqslant 2\pi$）的弧长、所围图形的面积以及绕 $Ox$ 轴旋转得到的旋转体的体积.}
}

% 38 = 例15.11
\vbox{%
\examp{38.}{求微分方程 $y'+1=\mathrm{e}^{-y}\sin x$ 的通解.}
}

% 39 = 例14.5
\vbox{%
\examp{39.}{设 $f(x,y)$ 在 $D=\left\{(x,y)\mid x^{2}+y^{2}\leqslant 1\right\}$ 上连续，$f(x,y)=\mathrm{e}^{x^{2}+y^{2}}-\displaystyle\iint\limits_{D}\dfrac{(2x^{2}+1)f(x,y)}{x^{2}+y^{2}+1}\,\mathrm{d}x\mathrm{d}y$，求 $\displaystyle\iint\limits_{D}f(x,y)\,\mathrm{d}\sigma$.}
}

% 40 = 例9.12
\vbox{%
\examp{40.}{求不定积分 $\displaystyle\int\sin^2 x\cos^3 x\,\mathrm{d}x$.}
}

% 41 = 例15.4
\vbox{%
\examp{41.}{设函数 $y(x)$ 是微分方程 $2xy'-4y=2\ln x-1$ 满足条件 $y(1)=\dfrac{1}{4}$ 的解，求曲线 $y=y(x)$（$1\leqslant x\leqslant \mathrm{e}$）的弧长.}
}

% 42 = 例11.13
\vbox{%
\examp{42.}{设 $I=\displaystyle\int_{0}^{\frac{3\pi}{2}}\dfrac{\cos x}{2x-3\pi}\,\mathrm{d}x$，则 $I$（\quad）.}
\begin{opts4}
(A) 为正 & (B) 为负 & (C) 为零 & (D) 发散
\end{opts4}
}

% 43 = 例13.17
\vbox{%
\examp{43.}{设 $y=y(x)$ 及 $z=z(x)$ 由方程 $\mathrm{e}^z-xyz=0$ 及 $xz^2=\ln y$ 所确定，则 $\left.\dfrac{\mathrm{d}z}{\mathrm{d}x}\right|_{x=\frac{1}{2}}=\underline{\hspace{3em}}$.}
}

% 44 = 例5.3
\vbox{%
\examp{44.}{已知曲线 $y=f(x)$ 在点 $(0,1)$ 处的曲率圆方程为 $(x-1)^{2}+y^{2}=2$，且当 $x\to 0$ 时，二阶可导函数 $f(x)$ 与 $a+bx+cx^{2}$ 的差为 $o(x^{2})$，则（\quad）.}
\begin{opts2}
(A) $a=0$，$b=1$，$c=\frac{3}{2}$ & (B) $a=1$，$b=0$，$c=1$ \\
(C) $a=1$，$b=1$，$c=-1$ & (D) $a=1$，$b=0$，$c=-1$
\end{opts2}
}

% 45 = 例3.3
\vbox{%
\examp{45.}{设函数 $f(x)$ 连续，给出下列 $4$ 个条件：\\
(1) $\lim\limits_{x\to 0}\dfrac{|f(x)|-f(0)}{x}$ 存在；\\
(2) $\lim\limits_{x\to 0}\dfrac{f(x)-|f(0)|}{x}$ 存在；\\
(3) $\lim\limits_{x\to 0}\dfrac{|f(x)|}{x}$ 存在；\\
(4) $\lim\limits_{x\to 0}\dfrac{|f(x)|-|f(0)|}{x}$ 存在.\\
其中可得到“$f(x)$ 在 $x=0$ 处可导”的条件个数为（\quad）.}
\begin{opts4}
(A) $1$ & (B) $2$ & (C) $3$ & (D) $4$
\end{opts4}
}

% 46 = 例8.4
\vbox{%
\examp{46.}{求极限 $\lim\limits_{n\to\infty}\sum\limits_{k=1}^{n}\dfrac{\pi}{4n}\cos^{2}\dfrac{k\pi}{4n}$.}
}

% 47 = 例6.5
\vbox{%
\examp{47.}{证明：若 $f(x)$ 连续且满足 $\displaystyle\int_0^{\frac{\pi}{2}}f(x)\cos x\,\mathrm{d}x=0$，则存在 $\xi\in\left(0,\frac{\pi}{2}\right)$，使得 $f(\xi)=0$.}
}

% 48 = 例9.16
\vbox{%
\examp{48.}{设函数 $f(x)$ 在 $[0,+\infty)$ 内可导，$f(0)=0$，其反函数为 $g(x)$. 若 $\displaystyle\int_x^{x+f(x)}g(t-x)\,\mathrm{d}t=x^2\ln(1+x)$，求 $f(x)$.}
}

% 49 = 例9.4
\vbox{%
\examp{49.}{求 $\displaystyle\int\sqrt{\dfrac{2x+1}{x+1}}\,\mathrm{d}x$.}
}

% 50 = 例10.4
\vbox{%
\examp{50.}{求摆线 $\begin{cases}x=a(t-\sin t),\\ y=a(1-\cos t)\end{cases}$（$a>0$）一拱的弧长、与 $x$ 轴所围图形的面积及绕 $x$ 轴旋转一周得到的旋转体的体积.}
}

% 51 = 例11.7
\vbox{%
\examp{51.}{$I=\displaystyle\int_{0}^{1}\dfrac{\ln(1+x)}{1+x^{2}}\,\mathrm{d}x=\underline{\hspace{3em}}$.}
}

% 52 = 例13.21
\vbox{%
\examp{52.}{设 $z=z(u,v)$ 具有二阶连续偏导数，且 $z=z(x-2y,\,x+3y)$ 满足 $6\dfrac{\partial^2 z}{\partial x^2}+\dfrac{\partial^2 z}{\partial x\partial y}-\dfrac{\partial^2 z}{\partial y^2}=3\dfrac{\partial z}{\partial x}-\dfrac{\partial z}{\partial y}$，$\dfrac{\mathrm{d}[z(0,v)]}{\mathrm{d}v}=\dfrac{1}{5}z(0,v)+\mathrm{e}^v$，$z(u,0)=u^2$. \\
(1) 证明 $\dfrac{\partial^2 z}{\partial u\partial v}=\dfrac{1}{5}\cdot\dfrac{\partial z}{\partial u}$；\\
(2) 求 $z=z(u,v)$ 的表达式.}
}

% 53 = 例1.15
\vbox{%
\examp{53.}{已知函数 $f(x)$ 的定义域是 $[0,+\infty)$，且满足 $f(0)=1$，$f'(x)=\dfrac{1}{f^{2}(x)+x^{2}}$，求证：$\lim\limits_{x\to+\infty}f(x)$ 存在，且 $\lim\limits_{x\to+\infty}f(x)\leqslant 1+\dfrac{\pi}{2}$.}
}

% 54 = 例2.1
\vbox{%
\examp{54.}{若 $x_{1}=1$，$x_{n+1}=\sqrt{4+3x_{n}}$，$n=1,2,\cdots$，证明数列 $\{x_{n}\}$ 收敛，并求 $\lim\limits_{n\to\infty}x_{n}$.}
}

% 55 = 例15.1
\vbox{%
\examp{55.}{微分方程 $x+yy'=y-xy'$ 的通解为\underline{\hspace{3em}}.}
}

% 56 = 例9.5
\vbox{%
\examp{56.}{设 $a_n=\displaystyle\int_0^{n\pi}x|\sin x|\,\mathrm{d}x$，$n=1,2,\cdots$，求 $a_n$ 的表达式.}
}

% 57 = 例9.14
\vbox{%
\examp{57.}{求不定积分 $\displaystyle\int\dfrac{1}{1+2\cos x}\,\mathrm{d}x$.}
}

% 58 = 例6.16
\vbox{%
\examp{58.}{设函数 $f(x)$ 在 $[0,1]$ 上二阶可导，$f(0)=f(1)$，且 $|f''(x)|\leqslant 2$，证明：$|f'(x)|\leqslant 1$，$x\in[0,1]$.}
}

% 59 = 例3.6
\vbox{%
\examp{59.}{设函数 $f(x)$ 处处可导，$f(0)=-1$，$f'(0)=1$，令 $g(x)=\bigl|f(x-1)\bigr|$，则 $g'(1)=\underline{\hspace{3em}}$.}
}

% 60 = 例1.12
\vbox{%
\examp{60.}{设 $f(x)$ 具有一阶连续导数，且 $\lim\limits_{x\to+\infty}[x-f(x)]=0$，$f(1)<1$，证明：\\
(1) 存在 $\xi>1$，使得 $|\xi-f(\xi)|<1-f(1)$；\\
(2) 存在 $\eta>1$，使得 $f'(\eta)>1$.}
}

% 61 = 例6.14
\vbox{%
\examp{61.}{设 $f(x)$ 在 $[a,b]$ 上连续，且 $f(x)>0$，证明存在 $\xi\in(a,b)$，使得 $\dfrac{b^{2}-a^{2}}{\displaystyle\int_{a}^{b}f(x)\,\mathrm{d}x}=\dfrac{2\xi}{f(\xi)}$.}
}

% 62 = 例3.4
\vbox{%
\examp{62.}{设 $f(x)$ 在 $x=0$ 处连续，下列结论：\\
(1) 若 $f'_-(0)$ 存在，则 $f'_+(0)$ 存在；\\
(2) 若 $\lim\limits_{x\to 0^{-}}f'(x)$ 存在，则 $\lim\limits_{x\to 0^{+}}f'(x)$ 存在；\\
(3) 若 $f'_-(0)$ 与 $f'_+(0)$ 均存在，则 $f'(0)$ 存在；\\
(4) 若 $\lim\limits_{x\to 0}f'(x)$ 存在，则 $f'(0)$ 存在.\\
正确结论的个数为（\quad）.}
\begin{opts4}
(A) $1$ & (B) $2$ & (C) $3$ & (D) $4$
\end{opts4}
}

% 63 = 例6.12
\vbox{%
\examp{63.}{设函数 $f(x)$ 在区间 $[0,2]$ 上具有连续导数，$f(0)=f(2)=0$，$M=\max\limits_{x\in[0,2]}\{|f(x)|\}$. 证明：\\
(1) 存在 $\xi\in(0,2)$，使得 $|f'(\xi)|\geqslant M$；\\
(2) 若对任意的 $x\in(0,2)$，$|f'(x)|\leqslant M$，则 $M=0$.}
}

% 64 = 例11.11
\vbox{%
\examp{64.}{计算下列积分：\\[0.6ex]
(1) $\displaystyle\int_{0}^{2}(x-1)\,\mathrm{d}x$；\\[0.6ex]
(2) $\displaystyle\int_{0}^{2}x(x-1)(x-2)\,\mathrm{d}x$.}
}

% 65 = 例2.3
\vbox{%
\examp{65.}{若 $0<x_n<1$，$x_{n+1}=1-\sqrt{1-x_n}$（$n=1,2,\cdots$），求：\\
(1) $\lim\limits_{n\to\infty}x_n$；\\
(2) $\lim\limits_{n\to\infty}\dfrac{x_{n+1}}{x_n}$.}
}

% 66 = 例5.4
\vbox{%
\examp{66.}{设函数 $f(x)$ 在 $x=x_{0}$ 处有二阶导数，则（\quad）.}
\begin{opts1}
(A) 当 $f(x)$ 在 $x_{0}$ 的某邻域内单调增加时，$f'(x_{0})>0$ \\
(B) 当 $f'(x_{0})>0$ 时，$f(x)$ 在 $x_{0}$ 的某邻域内单调增加 \\
(C) 当曲线 $f(x)$ 在 $x_{0}$ 的某邻域内是凹的时，$f''(x_{0})>0$ \\
(D) 当 $f''(x_{0})>0$ 时，曲线 $f(x)$ 在 $x_{0}$ 的某邻域内是凹的
\end{opts1}
}

% 67 = 例13.15
\vbox{%
\examp{67.}{设 $z=z(x,y)$ 是由方程 $\dfrac{x}{z}=\ln\dfrac{z}{y}$ 确定的二元隐函数，求全微分 $\mathrm{d}z$.}
}

% 68 = 例13.13
\vbox{%
\examp{68.}{设 $f(x,y)=\begin{cases}xy, & |x|\geqslant|y|,\\ -xy, & |x|<|y|,\end{cases}$ 则 $f''_{xy}(0,0)$ 和 $f''_{yx}(0,0)$ 依次为（\quad）.}
\begin{opts4}
(A) $1$，$1$ & (B) $1$，$-1$ & (C) $-1$，$1$ & (D) $-1$，$-1$
\end{opts4}
}

% 69 = 例8.8
\vbox{%
\examp{69.}{反常积分 (1) $\displaystyle\int_0^{+\infty}\dfrac{|\sin x|}{\sqrt{x^3}}\,\mathrm{d}x$，(2) $\displaystyle\int_0^{+\infty}\dfrac{1}{1+\sqrt{x}\,|\sin x|}\,\mathrm{d}x$ 的敛散性为（\quad）.}
\begin{opts2}
(A) (1)收敛，(2)收敛 & (B) (1)收敛，(2)发散 \\
(C) (1)发散，(2)收敛 & (D) (1)发散，(2)发散
\end{opts2}
}

% 70 = 例6.4
\vbox{%
\examp{70.}{设正值函数 $f(x)$，$g(x)$ 在 $[a,b]$ 上连续，证明存在 $\xi\in(a,b)$，使得 $\dfrac{f(\xi)}{g(\xi)}=\dfrac{\displaystyle\int_a^\xi f(x)\,\mathrm{d}x}{\displaystyle\int_\xi^b g(x)\,\mathrm{d}x}$.}
}

% 71 = 例13.20
\vbox{%
\examp{71.}{设函数 $z=xyf\left(\dfrac{y}{x}\right)$，其中 $f(u)$ 可导，若 $x\dfrac{\partial z}{\partial x}+y\dfrac{\partial z}{\partial y}=y^2(\ln y-\ln x)$，则（\quad）.}
\begin{opts2}
(A) $f(1)=\frac{1}{2}$，$f'(1)=0$ & (B) $f(1)=0$，$f'(1)=1$ \\
(C) $f(1)=\frac{1}{2}$，$f'(1)=1$ & (D) $f(1)=0$，$f'(1)=\frac{1}{2}$
\end{opts2}
}

% 72 = 例13.27
\vbox{%
\examp{72.}{设 $u(x,y)$ 在平面有界闭区域 $D$ 上具有二阶连续偏导数，且 $\dfrac{\partial^2 u}{\partial x\partial y}\neq 0$，$\dfrac{\partial^2 u}{\partial x^2}\cdot\dfrac{\partial^2 u}{\partial y^2}=0$，则 $u(x,y)$ 的（\quad）.}
\begin{opts1}
(A) 最大值点和最小值点必定都在 $D$ 的内部 \\
(B) 最大值点和最小值点必定都在 $D$ 的边界上 \\
(C) 最大值点在 $D$ 的内部，最小值点在 $D$ 的边界上 \\
(D) 最小值点在 $D$ 的内部，最大值点在 $D$ 的边界上
\end{opts1}
}

% 73 = 例8.6
\vbox{%
\examp{73.}{设常数 $p>0$，$q>0$，若 $\displaystyle\int_{0}^{1}\dfrac{\ln x}{x^{p}(1-x)^{q}}\,\mathrm{d}x$ 收敛，则（\quad）.}
\begin{opts2}
(A) $0<p<1,\ 0<q<2$ & (B) $p>1,\ 1<q<2$ \\
(C) $0<p<1,\ 0<q<1$ & (D) $p>1,\ 0<q<1$
\end{opts2}
}

% 74 = 例2.5
\vbox{%
\examp{74.}{设 $x_{n+1}=\sqrt{x_n(2-x_n)}$（$n=1,2,\cdots$），$0<x_1<2$，证明数列 $\{x_n\}$ 的极限存在，并求此极限.}
}

% 75 = 例5.2
\vbox{%
\examp{75.}{设函数 $f(x)=(x^{2}+a)\mathrm{e}^{x}$，若 $f(x)$ 没有极值点，但曲线 $y=f(x)$ 有拐点，则 $a$ 的取值范围是（\quad）.}
\begin{opts4}
(A) $[0,1)$ & (B) $[1,+\infty)$ & (C) $[1,2)$ & (D) $[2,+\infty)$
\end{opts4}
}

% 76 = 例8.5
\vbox{%
\examp{76.}{设 $f(x)$ 连续，则 $\lim\limits_{n\to\infty}\sum\limits_{k=1}^{n}\dfrac{k-\dfrac{1}{2}+n}{n^{2}}f\left(\dfrac{2k-1}{2n}\right)=$（\quad）.}
\begin{opts2}
(A) $\int_{0}^{1}\left(x+\frac{1}{2}\right)f(x)\,\mathrm{d}x$ & (B) $\int_{0}^{1}\left(x+1-\frac{1}{2n}\right)f\left(x-\frac{1}{2n}\right)\mathrm{d}x$ \\
(C) $\int_{0}^{1}(x+1)f(x)\,\mathrm{d}x$ & (D) $\int_{0}^{1}\left(x+\frac{1}{2}+\frac{1}{n}\right)f\left(x-\frac{1}{n}\right)\mathrm{d}x$
\end{opts2}
}

% 77 = 例11.14
\vbox{%
\examp{77.}{设 $f(x)=xg'(2x)$，且 $g(x)$ 的一个原函数为 $\ln(x+1)$，则 $\displaystyle\int_{0}^{1}f(x)\,\mathrm{d}x=\underline{\hspace{3em}}$.}
}

% 78 = 例4.5
\vbox{%
\examp{78.}{设函数 $y=f(x)$ 由 $\begin{cases} x=2t+|t|,\\ y=|t|\tan t \end{cases}$ 所确定，则在 $\left(-\dfrac{\pi}{2},\dfrac{\pi}{2}\right)$ 内（\quad）.}
\begin{opts2}
(A) $f(x)$ 连续，$f'(0)$ 不存在 & (B) $f'(0)$ 存在，$f'(x)$ 在 $x=0$ 处不连续 \\
(C) $f'(x)$ 连续，$f''(0)$ 不存在 & (D) $f''(0)$ 存在，$f''(x)$ 在 $x=0$ 处不连续
\end{opts2}
}

% 79 = 例12.3
\vbox{%
\examp{79.}{半径为 $a$ 的半球形水池蓄满了水，水的比重为 1，现将水抽干，至少做功 $\dfrac{\pi}{2}$，则 $a=\underline{\hspace{3em}}$.}
}

% 80 = 例10.12
\vbox{%
\examp{80.}{求曲线 $(y+1)^{2}=(2-x)\ln x$ 在区间 $[1,2]$ 上所围成的平面图形绕直线 $y=-1$ 旋转一周所生成的旋转体的体积.}
}

% 81 = 例5.9
\vbox{%
\examp{81.}{设 $f'(x)$ 在区间 $[0,4]$ 上连续，曲线 $y=f'(x)$ 与直线 $x=0$，$x=4$，$y=0$ 围成如图所示的三个区域，其面积分别为 $S_{1}=3$，$S_{2}=4$，$S_{3}=2$，且 $f(0)=1$，则 $f(x)$ 在 $[0,4]$ 上的最大值与最小值分别为（\quad）.}
\begin{opts4}
(A) $2,-3$ & (B) $4,-3$ & (C) $2,-2$ & (D) $4,-2$
\end{opts4}
}

% 82 = 例1.5
\vbox{%
\examp{82.}{设函数 $f(x)$，$g(x)$ 在 $x=0$ 的某去心邻域内有定义且恒不为 $0$，若当 $x\to 0$ 时，$f(x)$ 是 $g(x)$ 的高阶无穷小，则当 $x\to 0$ 时，有（\quad）}
\begin{opts2}
(A) $f(x)+g(x)=o(g(x))$ & (B) $f(x)g(x)=o(f^{2}(x))$ \\
(C) $f(x)=o(e^{g(x)}-1)$ & (D) $f(x)=o(g^{2}(x))$
\end{opts2}
}

% 83 = 例9.17
\vbox{%
\examp{83.}{设 $|x|\leqslant 1$，求函数 $f(x)=\displaystyle\int_{-1}^{1}|t-x|\,\mathrm{e}^{2t}\,\mathrm{d}t$ 的最大值.}
}

% 84 = 例6.6
\vbox{%
\examp{84.}{设函数 $f(x)$ 在 $[0,\pi]$ 上连续，且 $\displaystyle\int_0^\pi f(x)\,\mathrm{d}x=0$，$\displaystyle\int_0^\pi f(x)\cos x\,\mathrm{d}x=0$. 证明：在 $(0,\pi)$ 内至少存在两个不同的点 $\xi_1$，$\xi_2$，使得 $f(\xi_1)=f(\xi_2)=0$.}
}

% 85 = 例9.2
\vbox{%
\examp{85.}{求下列不定积分：\\[0.6ex]
(1) $\displaystyle\int\sqrt{a^2-x^2}\,\mathrm{d}x$（$a>0$）；\\[0.6ex]
(2) $\displaystyle\int\sqrt{a^2+x^2}\,\mathrm{d}x$（$a>0$）；\\[0.6ex]
(3) $\displaystyle\int\sqrt{x^2-a^2}\,\mathrm{d}x$（$a>0$）.}
}

% 86 = 例15.17
\vbox{%
\examp{86.}{一质量为 $m$ 的飞机，着陆时的水平速度为 $v(0)$. 经测试，减速伞打开后，飞机所受的总阻力与飞机的速度成正比，比例系数为 $k$（$k>0$）. 从飞机接触跑道开始计时，设 $t$ 时刻飞机的滑行距离为 $x(t)$，速度为 $v(t)$，则飞机在滑行至总位移的一半时的速度为（\quad）.}
\begin{opts4}
(A) $\dfrac{v(0)}{3}$ & (B) $\dfrac{v(0)}{2}$ & (C) $\dfrac{v(0)}{\sqrt{3}}$ & (D) $\dfrac{v(0)}{\sqrt{2}}$
\end{opts4}
}

% 87 = 例13.22
\vbox{%
\examp{87.}{设 $z=z(x,y)$ 有二阶连续偏导数，用变换 $u=x-2y$，$v=x+ay$ 可把方程 $6\dfrac{\partial^2 z}{\partial x^2}+\dfrac{\partial^2 z}{\partial x\partial y}-\dfrac{\partial^2 z}{\partial y^2}=0$ 化简为 $\dfrac{\partial^2 z}{\partial u\partial v}=0$，求常数 $a$.}
}

% 88 = 例14.4
\vbox{%
\examp{88.}{确定积分区域 $D$，使得二重积分 $I=\displaystyle\iint\limits_{D}\left(1-x^{2}-\dfrac{y^{2}}{2}\right)\mathrm{d}x\mathrm{d}y$ 达到最大值.}
}

% 89 = 例5.6
\vbox{%
\examp{89.}{曲线 $y=\dfrac{(1+x)^{\frac{3}{2}}}{\sqrt{x}}$ 的斜渐近线方程为 $\underline{\hspace{3em}}$.}
}

% 90 = 例5.5
\vbox{%
\examp{90.}{曲线 $y=x\ln\left(\mathrm{e}+\dfrac{1}{x-1}\right)$ 的斜渐近线方程为（\quad）.}
\begin{opts4}
(A) $y=x+\mathrm{e}$ & (B) $y=x+\frac{1}{\mathrm{e}}$ & (C) $y=x$ & (D) $y=x-\frac{1}{\mathrm{e}}$
\end{opts4}
}

% 91 = 例12.2
\vbox{%
\examp{91.}{设某物体运动的路程 $s$ 和加速度 $a$ 与时间 $t$ 的关系分别为 $s=s(t)$，$a=a(t)$. 已知该物体从 $t=0$ 时刻至 $t=1$ 时刻运动的路程是 2，在 $t=1$ 时刻的速度为 2，则 $\displaystyle\int_0^1 t\,a(t)\,\mathrm{d}t=$（\quad）.}
\begin{opts4}
(A) $0$ & (B) $1$ & (C) $2$ & (D) $3$
\end{opts4}
}

% 92 = 例10.15
\vbox{%
\examp{92.}{已知 $M(x_{0},y_{0})$ 是曲线 $y=\dfrac{1}{1+x^{2}}$（$x\geqslant 0$）的拐点，$O$ 为坐标原点. 记 $D$ 是第一象限中以曲线 $y=\dfrac{1}{1+x^{2}}$（$x\geqslant x_{0}$）、线段 $OM$ 及 $x$ 正半轴为边界的无界区域，求 $D$ 绕 $x$ 轴旋转所成旋转体的体积.}
}

% 93 = 例15.9
\vbox{%
\examp{93.}{欧拉方程 $x^{2}y''+xy'-4y=0$ 满足条件 $y(1)=1$，$y'(1)=2$ 的解为 $y=\underline{\hspace{3em}}$.}
}

% 94 = 例10.8
\vbox{%
\examp{94.}{设函数 $f(x)$ 满足 $\displaystyle\int\dfrac{f(x)}{\sqrt{x}}\,\mathrm{d}x=\dfrac{1}{6}x^2-x+C$，$L$ 为曲线 $y=f(x)$（$4\leqslant x\leqslant 9$）. 求 $L$ 绕 $x$ 轴旋转所成旋转曲面的面积.}
}

% 95 = 例14.7
\vbox{%
\examp{95.}{设平面有界区域 $D$ 位于第一象限，由曲线 $x^{2}+y^{2}-xy=1$，$x^{2}+y^{2}-xy=2$ 与直线 $y=\sqrt{3}x$，$y=0$ 围成，计算 $\displaystyle\iint\limits_{D}\dfrac{1}{3x^{2}+y^{2}}\,\mathrm{d}x\mathrm{d}y$.}
}

% 96 = 例9.21
\vbox{%
\examp{96.}{设 $f(x)=\begin{cases}\mathrm{e}^{-x}, & x\geqslant 0,\\ 1+x^2, & x<0\end{cases}$，则 $\displaystyle\int_{-2}^{2}f(x-1)\,\mathrm{d}x=\underline{\hspace{3em}}$.}
}

% 97 = 例13.9
\vbox{%
\examp{97.}{函数 $g(x,y)=\sqrt{x^2+y^2}$ 在点 $(0,0)$ 处（\quad）.}
\begin{opts2}
(A) 连续，偏导数存在 & (B) 连续，偏导数不存在 \\
(C) 不连续，偏导数存在 & (D) 不连续，偏导数不存在
\end{opts2}
}

% 98 = 例12.1
\vbox{%
\examp{98.}{如图所示，井深 $a$ 米，每米绳子的重量是 $5\,\mathrm{N}$，挂斗重 $400\,\mathrm{N}$，污泥重 $1500\,\mathrm{N}$. 将挂斗从井底提到井口所做的功为 $59250\,\mathrm{J}$，则 $a=\underline{\hspace{3em}}$.}
}

% 99 = 例15.12
\vbox{%
\examp{99.}{设 $x\geqslant -2$，则微分方程 $y'-\sqrt{y+x^{2}}=-2x$ 满足 $y(0)=1$ 的特解为\underline{\hspace{3em}}.}
}

% 100 = 例5.7
\vbox{%
\examp{100.}{设 $g(x)$ 是函数 $f(x)=\dfrac{1}{2}\ln\dfrac{3+x}{3-x}$ 的反函数，则曲线 $y=g(x)$ 的渐近线方程为 $\underline{\hspace{3em}}$.}
}

% 101 = 例3.7
\vbox{%
\examp{101.}{下列函数中，在 $x=0$ 处不可导的是（\quad）.}
\begin{opts2}
(A) $f(x)=|x|\tan|x|$ & (B) $f(x)=|x|\tan\sqrt{|x|}$ \\
(C) $f(x)=\sqrt{\cos|x|}$ & (D) $f(x)=\cos\sqrt{|x|}$
\end{opts2}
}

% 102 = 例10.2
\vbox{%
\examp{102.}{求阿基米德螺线 $r=a\theta$（$a>0$，$0\leqslant\theta\leqslant 2\pi$）与 $Ox$ 轴所围图形的面积.}
}

% 103 = 例6.20
\vbox{%
\examp{103.}{已知常数 $k\geqslant\ln 2-1$. 证明：$(x-1)\left(x-\ln^2 x+2k\ln x-1\right)\geqslant 0$.}
}

% 104 = 例13.7
\vbox{%
\examp{104.}{函数 $f(x,y)=\begin{cases}1, & xy=0,\\ 0, & \text{其他},\end{cases}$ 在点 $(0,0)$ 处（\quad）.}
\begin{opts1}
(A) 关于两个变量都连续，其在原点连续 \\
(B) 关于两个变量都连续，其在原点不连续 \\
(C) 关于两个变量都不连续，其在原点连续 \\
(D) 关于两个变量都不连续，其在原点不连续
\end{opts1}
}

% 105 = 例15.13
\vbox{%
\examp{105.}{已知 $f(xy)=yf(x)+xf(y)$ 对任意正实数 $x,y$ 均成立，且 $f'(1)=\mathrm{e}$，求 $f(xy)$ 的极小值.}
}

% 106 = 例9.9
\vbox{%
\examp{106.}{设 $n$ 为非负整数，则 $\displaystyle\int_0^1 x^2\ln^n x\,\mathrm{d}x=\underline{\hspace{3em}}$.}
}

% 107 = 例14.2
\vbox{%
\examp{107.}{$\displaystyle\int_{0}^{\frac{\pi}{6}}\mathrm{d}y\displaystyle\int_{y}^{\frac{\pi}{6}}\dfrac{\cos x}{x}\,\mathrm{d}x=\underline{\hspace{3em}}$.}
}

% 108 = 例11.15
\vbox{%
\examp{108.}{设 $f(x)=\displaystyle\int_{0}^{x}\mathrm{e}^{-t^{2}+2t}\,\mathrm{d}t$，则 $\displaystyle\int_{0}^{1}(x-1)^{2}f(x)\,\mathrm{d}x=\underline{\hspace{3em}}$.}
}

% 109 = 例13.31
\vbox{%
\examp{109.}{求函数 $f(x,y)=3(x^2+y^2)-x^3$ 在有界闭区域 $D=\{(x,y)\mid x^2+y^2\leqslant 16\}$ 上的最大值和最小值.}
}

% 110 = 例15.18
\vbox{%
\examp{110.}{在 $xOy$ 平面上，设 $|PQ|=1$，初始时刻 $P$ 在原点，$Q$ 在 $(1,0)$ 点，若 $P$ 点沿着 $y$ 轴的正方向移动，且 $Q$ 点的运动方向始终指向 $P$ 点，求 $Q$ 点的运动轨迹.}
}

% 111 = 例11.5
\vbox{%
\examp{111.}{$I=\displaystyle\int_{-\frac{\pi}{2}}^{\frac{\pi}{2}}\dfrac{\sin^{4}x}{1+\mathrm{e}^{-x}}\,\mathrm{d}x=\underline{\hspace{3em}}$.}
}

% 112 = 例13.10
\vbox{%
\examp{112.}{函数 $f(x,y)=\begin{cases}\dfrac{x^2y}{x^2+y^2}, & (x,y)\neq(0,0),\\ 0, & (x,y)=(0,0)\end{cases}$ 在点 $(0,0)$ 处（\quad）.}
\begin{opts4}
(A) 连续，可微 & (B) 连续，不可微 & (C) 不连续，可微 & (D) 不连续，不可微
\end{opts4}
}

% 113 = 例4.1
\vbox{%
\examp{113.}{设函数 $f(x)$ 可导且满足 $x^{2}f'(x)=f^{2}(x)$，$f(1)=\dfrac{1}{3}$，则 $f^{(n)}(0)=$（\quad）.}
\begin{opts4}
(A) $(-1)^{n}n!$ & (B) $(-1)^{n-1}n!$ & (C) $(-2)^{n}n!$ & (D) $(-2)^{n-1}n!$
\end{opts4}
}

% 114 = 例8.1
\vbox{%
\examp{114.}{求极限 $\lim\limits_{n\to\infty}\dfrac{1}{n}\sqrt[n]{n(n+1)(n+2)\cdots\left(n+n-1\right)}$.}
}

% 115 = 例9.6
\vbox{%
\examp{115.}{求 $\displaystyle\int\dfrac{x\mathrm{e}^x}{\sqrt{\mathrm{e}^x+1}}\,\mathrm{d}x$.}
}

% 116 = 例9.3
\vbox{%
\examp{116.}{求下列不定积分：\\[0.6ex]
(1) $\displaystyle\int\dfrac{1}{x^2\sqrt{1+x^2}}\,\mathrm{d}x$；\\[0.6ex]
(2) $\displaystyle\int\dfrac{1}{\sqrt{\mathrm{e}^x+1}}\,\mathrm{d}x$；\\[0.6ex]
(3) $\displaystyle\int x\sqrt{x^2+2x+2}\,\mathrm{d}x$.}
}

% 117 = 例10.9
\vbox{%
\examp{117.}{已知 $y=\sqrt{x}\,a^{\frac{x}{2a}}$（$a>1$，$0\leqslant x<+\infty$）与 $x$ 轴中间部分区域绕 $x$ 轴旋转一周生成的旋转体的体积为 $\pi\mathrm{e}^2$，则 $a=\underline{\hspace{3em}}$.}
}

% 118 = 例1.6
\vbox{%
\examp{118.}{已知函数 $f(x)=\dfrac{\displaystyle\int_{0}^{x}\ln(1+t^{2})\,\mathrm{d}t}{x^{\alpha}}$ 在 $(0,+\infty)$ 上有界，则 $\alpha$ 的取值范围是（\quad）.}
\begin{opts4}
(A) $-1<\alpha<3$ & (B) $1<\alpha\leqslant 3$ & (C) $0<\alpha\leqslant 3$ & (D) $\alpha>3$
\end{opts4}
}

% 119 = 例15.2
\vbox{%
\examp{119.}{微分方程 $\dfrac{\mathrm{d}y}{\mathrm{d}x}=\dfrac{2x-5y+3}{2x+4y-6}$ 满足 $y(0)=2$ 的特解为\underline{\hspace{3em}}.}
}

% 120 = 例11.6
\vbox{%
\examp{120.}{(1) 设 $I=\displaystyle\int_{a}^{a+\pi}|\sin nx|\,\mathrm{d}x$，$n=1,2,\cdots$，$a$ 为任意常数，则（\quad）.}
\begin{opts4}
(A) $I$ 只与 $a$ 有关 & (B) $I$ 只与 $n$ 有关 & (C) $I$ 与 $a$，$n$ 均有关 & (D) $I$ 与 $a$，$n$ 均无关
\end{opts4}
\noindent (2) 设 $I=\displaystyle\int_{a}^{a+\frac{k\pi}{2}}\sqrt{1-\sin^{2}x}\,\mathrm{d}x$，$k$ 为正整数，$a$ 为任意实数，则（\quad）.
}

% 121 = 例1.13
\vbox{%
\examp{121.}{设 $f(x)$ 单调减少，$\lim\limits_{x\to+\infty}f(x)=0$，证明 $f(x)\geqslant 0$.}
}

% 122 = 例3.5
\vbox{%
\examp{122.}{设 $f(x)=x^{\frac{2}{3}}\sin x$，求 $f'(x)$.}
}

% 123 = 例10.3
\vbox{%
\examp{123.}{求三叶玫瑰线 $r=a\sin 3\theta$（$a>0$）所围图形的面积.}
}

% 124 = 例10.5
\vbox{%
\examp{124.}{求星形线 $x^{\frac{2}{3}}+y^{\frac{2}{3}}=a^{\frac{2}{3}}$（$a>0$）的弧长、所围图形的面积及绕 $x$ 轴旋转一周所得旋转体的体积.}
}

% 125 = 例3.2
\vbox{%
\examp{125.}{设函数 $f(x)$ 在区间 $(-1,1)$ 内有定义，且在点 $x=0$ 处连续，则下列命题中：\\
(1) 当 $\lim\limits_{x\to 0}\dfrac{f(x)}{\sqrt[3]{x}}=0$ 时，$f(x)$ 在点 $x=0$ 处可导；\\
(2) 当 $\lim\limits_{x\to 0}\dfrac{f(x)}{x^{2}}=0$ 时，$f(x)$ 在点 $x=0$ 处可导；\\
(3) 当 $f(x)$ 在点 $x=0$ 处可导时，$\lim\limits_{x\to 0}\dfrac{f(x)}{\sqrt[3]{x}}=0$；\\
(4) 当 $f(x)$ 在点 $x=0$ 处可导时，$\lim\limits_{x\to 0}\dfrac{f(x)}{x^{2}}=0$.\\
真命题的个数为（\quad）.}
\begin{opts4}
(A) $1$ & (B) $2$ & (C) $3$ & (D) $4$
\end{opts4}
}

% 126 = 例6.1
\vbox{%
\examp{126.}{设函数 $f(x)$ 在 $[0,1]$ 上连续，$(0,1)$ 内可导，$f(0)=f(1)$，$\displaystyle\int_0^1 f(x)\,\mathrm{d}x=0$，证明：存在 $\xi\in(0,1)$，使得 $f'(\xi)+f^2(\xi)=0$.}
}

% 127 = 例15.6
\vbox{%
\examp{127.}{求微分方程 $2y^{2}y''=[1+(y')^{2}]^{2}$ 满足初始条件 $y(3)=2$，$y'(3)=1$ 的解.}
}

% 128 = 例11.18
\vbox{%
\examp{128.}{设函数 $f(x)$ 具有二阶导数，且 $f'(0)=f'(1)$，$|f''(x)|\leqslant 1$. 证明：\\
(1) 当 $x\in(0,1)$ 时，$|f(x)-f(0)(1-x)-f(1)x|\leqslant \dfrac{x(1-x)}{2}$；\\[0.6ex]
(2) $\left|\displaystyle\int_0^1 f(x)\,\mathrm{d}x - \dfrac{f(0)+f(1)}{2}\right|\leqslant \dfrac{1}{12}$.}
}

% 129 = 例10.13
\vbox{%
\examp{129.}{$y=\sqrt{x}\,\mathrm{e}^{-\frac{3}{2}x}$（$x\geqslant 0$）下方及 $x$ 轴上方的无界区域绕 $x$ 轴旋转一周所生成的旋转体的体积为\underline{\hspace{3em}}.}
}

% 130 = 例12.5
\vbox{%
\examp{130.}{设沿 $y$ 轴上的区间 $[0,1]$ 放置一长度为 1 且线密度为 $\rho$ 的均匀细杆，在 $x$ 轴上 $x=1$ 处有一单位质点，则该细杆对此质点的引力（$G$ 为引力常量）沿 $x$ 轴正向的分力为\underline{\hspace{3em}}.}
}

% 131 = 例3.1
\vbox{%
\examp{131.}{设 $f(x)=2-x^{x}-x+o(x-1)$，$x\to 1$，且 $f'(1)=a$，则 $a=\underline{\hspace{3em}}$.}
}

% 132 = 例13.1
\vbox{%
\examp{132.}{求 $\lim\limits_{(x,y)\to(0,0)}\dfrac{x^2-y^2}{x^2+y^2}$.}
}

% 133 = 例2.9
\vbox{%
\examp{133.}{设数列 $\{a_n\}$，$\{b_n\}$ 满足 $a_n+b_n=\dfrac{\pi}{2}$，$a_{n+1}=a_n+\cos a_n$，其中 $n=1,2,\cdots$，$a_1=1$，当 $n\to\infty$ 时，$b_n$ 与 $b_{n-1}^{a}$ 为同阶无穷小量，则 $a=(\quad)$.}
\begin{opts4}
(A) $1$ & (B) $2$ & (C) $3$ & (D) $5$
\end{opts4}
}

% 134 = 例1.4
\vbox{%
\examp{134.}{当 $x\to 0$ 时，$\cos\left(e^{\frac{x^{2}}{2}}-1\right)-1\sim cx^{k}$，则 $ck=\underline{\hspace{3em}}$.}
}

% 135 = 例6.9
\vbox{%
\examp{135.}{已知函数 $f(x)$ 在 $[x_{0},\,x_{0}+\delta)$ 上连续，在 $(x_{0},\,x_{0}+\delta)$ 内可导，$\delta>0$，证明：若 $\lim\limits_{x\to x_{0}^{+}}f'(x)=A$，则 $f'_{+}(x_{0})=A$.}
}

% 136 = 例13.16
\vbox{%
\examp{136.}{已知函数 $f(u,v)$ 具有二阶连续偏导数，$f(1,1)=2$ 是 $f(u,v)$ 的极值，$z=f[x+y,\,f(x,y)]$，求 $\left.\dfrac{\partial^2 z}{\partial x\partial y}\right|_{(1,1)}$.}
}

% 137 = 例10.6
\vbox{%
\examp{137.}{求笛卡儿叶形线 $x^3+y^3-3axy=0$（$a>0$，参数方程为 $\begin{cases}x=\dfrac{3at}{1+t^3},\\ y=\dfrac{3at^2}{1+t^3}\end{cases}$）所围图形的面积及渐近线.}
}

% 138 = 例8.7
\vbox{%
\examp{138.}{设 $p$ 为常数，若反常积分 $\displaystyle\int_0^1 x^p(1-x)^{p-1}\ln x\,\mathrm{d}x$ 收敛，则（\quad）.}
\begin{opts4}
(A) $p<-1$ & (B) $-1<p\leqslant 0$ & (C) $0\leqslant p<1$ & (D) $p>1$
\end{opts4}
}

% 139 = 例6.18
\vbox{%
\examp{139.}{设 $f(x)$ 在 $\left[0,\dfrac{3}{2}\pi\right]$ 上连续，在 $\left(0,\dfrac{3}{2}\pi\right)$ 内是函数 $\dfrac{\sin x}{x}$ 的一个原函数，$f(0)=0$. \\
(1) 证明 $f\left(\dfrac{3}{2}\pi\right)>0$；\\[0.6ex]
(2) 求方程 $\displaystyle\int_0^{x}\dfrac{\sin t}{t}\,\mathrm{d}t=\ln x^2$ 的实根个数.}
}

% 140 = 例6.3
\vbox{%
\examp{140.}{设 $f(x)$ 在 $[0,1]$ 上连续，在 $(0,1)$ 内可导，且满足 $f(1)=k\displaystyle\int_0^{\frac{1}{k}}x\mathrm{e}^{1-x}f(x)\,\mathrm{d}x\ (k>1)$，证明至少存在一点 $\xi\in(0,1)$，使得 $f'(\xi)=\left(1-\xi^{-1}\right)f(\xi)$.}
}

% 141 = 例6.22
\vbox{%
\examp{141.}{设可导函数 $f(x)$ 严格单调递增且满足 $\displaystyle\int_{-1}^{1}f(x)\,\mathrm{d}x=0$，记 $a=\displaystyle\int_0^{1}f(x)\,\mathrm{d}x$. \\[0.6ex]
(1) 证明 $a>0$；\\[0.6ex]
(2) 令 $F(x)=a(1-x^2)+\displaystyle\int_1^{x}f(t)\,\mathrm{d}t$，证明：存在 $\xi\in(-1,1)$，使得 $F''(\xi)=0$.}
}

% 142 = 例14.8
\vbox{%
\examp{142.}{设 $D=\left\{(x,y)\mid(x-1)^{2}+(y-1)^{2}\leqslant 2,\ y\geqslant x\right\}$，计算二重积分 $\displaystyle\iint\limits_{D}(x-y)\,\mathrm{d}x\mathrm{d}y$.}
}

% 143 = 例1.7
\vbox{%
\examp{143.}{$\lim\limits_{n\to\infty}n^{2}\bigl(\sqrt[n]{2}-\sqrt[n+1]{2}\bigr)=\underline{\hspace{5em}}$.}
}

% 144 = 例2.4
\vbox{%
\examp{144.}{设 $x_{n+1}=\dfrac{x_n^{2}}{2(x_n-1)}$（$n=1,2,\cdots$），$x_1>1$，证明数列 $\{x_n\}$ 收敛，并求 $\lim\limits_{n\to\infty}x_n$.}
}

% 145 = 例2.7
\vbox{%
\examp{145.}{若 $(1-x_{n+1})x_n>\dfrac{1}{4}$（$n=1,2,\cdots$），$0<x_n<1$，证明数列 $\{x_n\}$ 收敛，并求 $\lim\limits_{n\to\infty}x_n$.}
}

% 146 = 例13.29
\vbox{%
\examp{146.}{设 $f(x,y)$ 与 $g(x,y)$ 均为可微函数，且 $g'_y(x,y)\neq 0$. 已知 $(x_0,y_0)$ 是 $f(x,y)$ 在约束条件 $g(x,y)=0$ 下的一个极值点，下列选项正确的是（\quad）.}
\begin{opts2}
(A) 若 $f'_x(x_0,y_0)=0$，则 $f'_y(x_0,y_0)=0$ & (B) 若 $f'_x(x_0,y_0)=0$，则 $f'_y(x_0,y_0)\neq 0$ \\
(C) 若 $f'_x(x_0,y_0)\neq 0$，则 $f'_y(x_0,y_0)=0$ & (D) 若 $f'_x(x_0,y_0)\neq 0$，则 $f'_y(x_0,y_0)\neq 0$
\end{opts2}
}

% 147 = 例11.2
\vbox{%
\examp{147.}{已知函数 $f(x)=\mathrm{e}^{\sin x}+\mathrm{e}^{-\sin x}$，则 $f'''(2\pi)=\underline{\hspace{3em}}$.}
}

% 148 = 例2.6
\vbox{%
\examp{148.}{(1) 证明对任意正整数 $n$，都有 $\dfrac{1}{n+1}<\ln\left(1+\dfrac{1}{n}\right)<\dfrac{1}{n}$ 成立； \\
(2) 设 $a_n=1+\dfrac{1}{2}+\cdots+\dfrac{1}{n}-\ln n$（$n=1,2,\cdots$），证明数列 $\{a_n\}$ 收敛.}
}

% 149 = 例8.3
\vbox{%
\examp{149.}{设 $f(x)=\begin{cases}\lim\limits_{n\to\infty}\dfrac{1}{n}\left(1+\cos\dfrac{x}{n}+\cos\dfrac{2x}{n}+\cdots+\cos\dfrac{n-1}{n}x\right), & x>0,\\ a, & x=0,\\ f(-x), & x<0\end{cases}$ 连续，则 $a=\underline{\hspace{3em}}$.}
}

% 150 = 例13.11
\vbox{%
\examp{150.}{函数 $f(x,y)=\begin{cases}x^2, & y=0,\\ y^2, & x=0,\\ 1, & \text{其他},\end{cases}$ 在点 $(0,0)$ 处（\quad）.}
\begin{opts1}
(A) 两个偏导数均连续，且函数可微 \\
(B) 两个偏导数均连续，且函数不可微 \\
(C) 两个偏导数均不连续，且函数可微 \\
(D) 两个偏导数均不连续，且函数不可微
\end{opts1}
}

% 151 = 例2.10
\vbox{%
\examp{151.}{设正项数列 $\{a_n\}$ 收敛于 $0$，若 $a_{n+1}=\cos b_n-\cos a_n$，$a_n\in\left(0,\dfrac{\pi}{2}\right)$，$b_n\in\left(0,\dfrac{\pi}{2}\right)$，且 $(1-b_n)^n=\cos b_n$，则 $\lim\limits_{n\to\infty}b_n^{\frac{\ln\cos b_n}{n}}=\underline{\hspace{3em}}$.}
}

% 152 = 例15.7
\vbox{%
\examp{152.}{求解微分方程 $xyy''+x(y')^{2}-yy'=0$.}
}

% 153 = 例2.13
\vbox{%
\examp{153.}{设正项数列 $\{a_n\}$ 单调增加，则以下选项中使得 $\{a_n\}$ 收敛的是（\quad）.}
\begin{opts2}
(A) $\left\{(1+a_n)^{\frac{1}{n}}\right\}$ 收敛于 $1$ & (B) $\left\{\left(1+\frac{1}{a_n}\right)^{a_n}\right\}$ 收敛于 $\mathrm{e}$ \\
(C) $\left\{a_n\ln a_n\right\}$ 收敛于 $0$ & (D) $\left\{\frac{\ln a_n}{a_n}\right\}$ 收敛于 $0$
\end{opts2}
}

% 154 = 例13.23
\vbox{%
\examp{154.}{设函数 $u=f\left(\ln\sqrt{x^2+y^2}\right)$ 有二阶连续偏导数，且满足 $\dfrac{\partial^2 u}{\partial x^2}+\dfrac{\partial^2 u}{\partial y^2}=(x^2+y^2)^{\frac{3}{2}}$，若极限 $\lim\limits_{x\to 0}\dfrac{\displaystyle\int_0^1 f(xt)\,\mathrm{d}t}{x}=-1$，求函数 $f(x)$ 的表达式.}
}

% 155 = 例13.12
\vbox{%
\examp{155.}{函数 $f(x,y)=\begin{cases}(x^2+y^2)\sin\dfrac{1}{x^2+y^2}, & (x,y)\neq(0,0),\\ 0, & (x,y)=(0,0),\end{cases}$ 在点 $(0,0)$ 处（\quad）.}
\begin{opts2}
(A) 可微，偏导数连续 & (B) 可微，但偏导数不连续 \\
(C) 不可微，偏导数连续 & (D) 不可微，且偏导数不连续
\end{opts2}
}

% 156 = 例15.10
\vbox{%
\examp{156.}{以 $y_{1}=t\mathrm{e}^{t}$，$y_{2}=\sin 2t$ 为两个特解的四阶常系数齐次线性微分方程为（\quad）.}
\begin{opts2}
(A) $y^{(4)}-2y'''+5y''-8y'+4y=0$ & (B) $y^{(4)}-2y'''+5y''+8y'+4y=0$ \\
(C) $y^{(4)}+2y'''+5y''-8y'+4y=0$ & (D) $y^{(4)}-2y'''-5y''-8y'+4y=0$
\end{opts2}
}

% 157 = 例15.8
\vbox{%
\examp{157.}{如果对微分方程 $y''-2ay'+(a+2)y=0$ 的任一解 $y(x)$，反常积分 $\displaystyle\int_{0}^{+\infty}y(x)\,\mathrm{d}x$ 均收敛，则 $a$ 的取值范围为（\quad）.}
\begin{opts4}
(A) $(-2,-1]$ & (B) $(-\infty,-1]$ & (C) $(-2,0)$ & (D) $(-\infty,0)$
\end{opts4}
}

% 158 = 例10.10
\vbox{%
\examp{158.}{曲线 $y=\displaystyle\int_{0}^{x}\sqrt{\cos t}\,\mathrm{d}t$ 的全长为\underline{\hspace{3em}}.}
}

% 159 = 例3.8
\vbox{%
\examp{159.}{已知函数 $g(x)$ 连续. 设 $f(x)=\displaystyle\int_{0}^{x^{2}}g(xt)\,\mathrm{d}t$，求 $f'(x)$ 的表达式，并判断 $f'(x)$ 在 $x=0$ 处的连续性.}
}

% 160 = 例14.6
\vbox{%
\examp{160.}{设 $D=\left\{(r,\theta)\mid 0\leqslant\theta\leqslant\dfrac{\pi}{4},\ 0\leqslant r\leqslant\sec\theta\right\}$，则 $\displaystyle\iint\limits_{D}\sqrt{1-(x-y)^{2}}\,\mathrm{d}\sigma=\underline{\hspace{3em}}$.}
}

% 161 = 例9.19
\vbox{%
\examp{161.}{极限 $\lim\limits_{t\to 0^+}\dfrac{1}{t^5}\displaystyle\int_0^t\mathrm{d}y\displaystyle\int_y^t\dfrac{\sin(xy)^2}{x}\,\mathrm{d}x=\underline{\hspace{3em}}$.}
}

% 162 = 例13.28
\vbox{%
\examp{162.}{设 $f(x)$ 为二阶可导函数，且 $x=0$ 是 $f(x)$ 的驻点，则二元函数 $z=f(x)f(y)$ 在点 $(0,0)$ 处取得极大值的一个充分条件是（\quad）.}
\begin{opts4}
(A) $f(0)<0$，$f''(0)>0$ & (B) $f(0)<0$，$f''(0)<0$ & (C) $f(0)>0$，$f''(0)>0$ & (D) $f(0)=0$，$f''(0)\neq 0$
\end{opts4}
}

% 163 = 例4.4
\vbox{%
\examp{163.}{设 $y=\dfrac{x}{1+x^{2}}$，求 $y^{(n)}$ 满足的递推关系式及 $y^{(2n+1)}(0)$.}
}

% 164 = 例6.11
\vbox{%
\examp{164.}{已知函数 $f(x)$ 在 $(-\infty,0)$ 上可导，且 $\lim\limits_{x\to-\infty}f'(x)=A>0$，证明 $\lim\limits_{x\to-\infty}f(x)=-\infty$.}
}

% 165 = 例9.13
\vbox{%
\examp{165.}{求不定积分 $\displaystyle\int\dfrac{2\sin x+\cos x}{\sin x+2\cos x}\,\mathrm{d}x$.}
}

% 166 = 例4.8
\vbox{%
\examp{166.}{若 $\begin{cases} x=\ln t,\\ y=\mathrm{e}^{t^{2}} \end{cases}$，则 $\left.\dfrac{\mathrm{d}^{2}y}{\mathrm{d}x^{2}}\right|_{t=1}=\underline{\hspace{3em}}$.}
}

% 167 = 例11.8
\vbox{%
\examp{167.}{$\displaystyle\int_{0}^{\pi}x\sqrt{\cos^{2}x-\cos^{4}x}\,\mathrm{d}x=\underline{\hspace{3em}}$.}
}

% 168 = 例6.8
\vbox{%
\examp{168.}{(1) 设 $f(x)$ 在 $[a,b]$ 上可导，若 $f'_{+}(a)\neq f'_{-}(b)$，证明：对于任意的介于 $f'_{+}(a)$ 与 $f'_{-}(b)$ 之间的 $\mu$，存在 $\xi\in(a,b)$，使得 $f'(\xi)=\mu$；\\
(2) 若 $f(x)$ 在 $[0,2]$ 上具有二阶导数，证明：存在 $\xi\in(0,2)$，使得 $f(0)+f(2)-2f(1)=f''(\xi)$.}
}

% 169 = 例13.25
\vbox{%
\examp{169.}{求函数 $f(x,y)=3(x^2+y^2)-x^3$ 的极值.}
}

% 170 = 例13.24
\vbox{%
\examp{170.}{设 $f(x,y)$ 是一阶偏导数连续的正值函数，满足 $f'_x(x,y)+f(x,y)=0$，若 $f'_y(0,y)=\tan y$，$f(0,0)=1$，求 $f(x,y)$.}
}

% 171 = 例11.12
\vbox{%
\examp{171.}{计算下列积分：\\[0.6ex]
(1) $\displaystyle\int_{0}^{4}x\sqrt{4x-x^{2}}\,\mathrm{d}x$；\\[0.6ex]
(2) $\displaystyle\int_{0}^{2}(2x+1)\sqrt{2x-x^{2}}\,\mathrm{d}x$.}
}

% 172 = 例15.3
\vbox{%
\examp{172.}{设一阶齐次线性微分方程 $y'+p(x)y=0$ 中，$p(x)$ 是以 $T$ 为周期的连续函数，则“该方程的非零解以 $T$ 为周期”是“$\displaystyle\int_{0}^{T}p(x)\,\mathrm{d}x=0$”的（\quad）.}
\begin{opts2}
(A) 充分非必要条件 & (B) 必要非充分条件 \\
(C) 充分必要条件 & (D) 既非充分也非必要条件
\end{opts2}
}

% 173 = 例10.14
\vbox{%
\examp{173.}{求曲线 $y=\dfrac{\mathrm{e}^{x}+\mathrm{e}^{-x}}{2}$，$x=0$，$x=t$（$t>0$），$y=0$ 所围成的曲边梯形绕 $x$ 轴旋转一周所得旋转体的侧面积.}
}

% 174 = 例6.17
\vbox{%
\examp{174.}{设函数 $f(x)=ax-b\ln x$（$a>0$）有 $2$ 个零点，则 $\dfrac{b}{a}$ 的取值范围是（\quad）.}
\begin{opts4}
(A) $(0,\mathrm{e})$ & (B) $(\mathrm{e},+\infty)$ & (C) $\left(0,\frac{1}{\mathrm{e}}\right)$ & (D) $\left(\frac{1}{\mathrm{e}},+\infty\right)$
\end{opts4}
}

% 175 = 例11.17
\vbox{%
\examp{175.}{设 $f(x)$ 在 $[0,1]$ 上连续，$\displaystyle\int_0^1 2x^2 f(x)\,\mathrm{d}x \geqslant \displaystyle\int_0^1 f^2(x)\,\mathrm{d}x + \dfrac{1}{5}$，则 $f(x)=\underline{\hspace{3em}}$.}
}

% 176 = 例13.26
\vbox{%
\examp{176.}{设 $f(x,y)=(y-x^2)(y-2x^2)$，$k$ 为任意常数，则（\quad）.}
\begin{opts1}
(A) $f(x,kx)$ 在 $x=0$ 处取极小值，点 $(0,0)$ 是 $f(x,y)$ 的极小值点 \\
(B) $f(x,kx)$ 在 $x=0$ 处取极小值，点 $(0,0)$ 不是 $f(x,y)$ 的极小值点 \\
(C) $f(x,kx)$ 在 $x=0$ 处不取极小值，点 $(0,0)$ 是 $f(x,y)$ 的极小值点 \\
(D) $f(x,kx)$ 在 $x=0$ 处不取极小值，点 $(0,0)$ 不是 $f(x,y)$ 的极小值点
\end{opts1}
}

% 177 = 例13.18
\vbox{%
\examp{177.}{设 $f(x,y)$ 具有一阶连续偏导数，证明：$f(x,y)$ 为 $k$ 次齐次函数（$f(tx,ty)=t^k f(x,y)$）的充要条件是 $x\dfrac{\partial f}{\partial x}+y\dfrac{\partial f}{\partial y}=kf(x,y)$.}
}

% 178 = 例12.6
\vbox{%
\examp{178.}{水从一根底面半径为 $1\,\mathrm{cm}$ 的圆柱形管道中流出. 因为水有黏性，在流动过程中受到管道壁的阻滞，所以流动的速度是随着到管道中心的距离而变化的，距管道中心越远，水流速度越小，在距离管道中心 $r\,\mathrm{cm}$ 处的水的流动速度为 $10(1-r^2)\,\mathrm{cm/s}$. 问水是以多大流量（以 $\mathrm{cm^3/s}$ 为单位）流过管道的？}
}

% 179 = 例11.10
\vbox{%
\examp{179.}{已知函数 $f(x)$ 在 $\left[0,\dfrac{\pi}{2}\right]$ 上可导，且 $\displaystyle\int_{0}^{\frac{\pi}{2}}f(x)\cos x\,\mathrm{d}x=0$，证明存在 $\xi\in\left(0,\dfrac{\pi}{2}\right)$，使得 $f'(\xi)=f(\xi)\tan\xi$.}
}

% 180 = 例10.7
\vbox{%
\examp{180.}{设 $y=\lim\limits_{n\to\infty}\dfrac{1+x}{1+nx^{2n}}$，则曲线 $y=y(x)$ 与 $x$ 轴及 $x=1$ 所围图形的面积为\underline{\hspace{3em}}.}
}

% 181 = 例1.14
\vbox{%
\examp{181.}{设 $x\geqslant 0$，记 $x$ 到 $2k$ 的最小距离为 $f(x)$，$k=0,1,2,\cdots$.\\
(1) 证明 $f(x)$ 以 2 为周期，写出其在 $[0,2]$ 上的表达式并画出 $f(x)$ 的图像；\\[0.6ex]
(2) 求 $\lim\limits_{x\to+\infty}\dfrac{\displaystyle\int_{0}^{x} f(t)\,\mathrm{d}t}{x}$.}
}

% 182 = 例1.8
\vbox{%
\examp{182.}{$\lim\limits_{x\to+\infty}x^{2}\left(3^{\frac{1}{x}}-3^{\frac{1}{x+1}}\right)=\underline{\hspace{5em}}$.}
}

% 183 = 例14.3
\vbox{%
\examp{183.}{$\displaystyle\int_{0}^{2}\mathrm{d}y\displaystyle\int_{y}^{2}\dfrac{y}{\sqrt{1+x^{3}}}\,\mathrm{d}x=\underline{\hspace{3em}}$.}
}

% 184 = 例4.3
\vbox{%
\examp{184.}{设 $f(x)=(x^{3}-1)^{n}$，则 $f^{(n)}(1)=\underline{\hspace{3em}}$.}
}

% 185 = 例11.3
\vbox{%
\examp{185.}{已知函数 $f(x)=\displaystyle\int_{0}^{x}\mathrm{e}^{\cos t}\,\mathrm{d}t$，$g(x)=\displaystyle\int_{0}^{\sin x}\mathrm{e}^{t^{2}}\,\mathrm{d}t$，则（\quad）.}
\begin{opts2}
(A) $f(x)$ 是奇函数，$g(x)$ 是偶函数 & (B) $f(x)$ 是偶函数，$g(x)$ 是奇函数 \\
(C) $f(x)$，$g(x)$ 均是奇函数 & (D) $f(x)$，$g(x)$ 均是周期函数
\end{opts2}
}

% 186 = 例13.19
\vbox{%
\examp{186.}{设 $z(x,y)=\sqrt{\dfrac{x^2+y^2}{xy}}$，$xy>0$，则 $x\dfrac{\partial z}{\partial x}+y\dfrac{\partial z}{\partial y}=\underline{\hspace{3em}}$.}
}

% 187 = 例13.2
\vbox{%
\examp{187.}{求 $\lim\limits_{(x,y)\to(0,0)}\dfrac{x^3+y^3}{x^2+y}$.}
}

% 188 = 例1.11
\vbox{%
\examp{188.}{函数 $f(x)=|x|^{\frac{1}{(1-x)(x-2)}}$ 的第一类间断点的个数是（\quad）.}
\begin{opts4}
(A) $3$ & (B) $2$ & (C) $1$ & (D) $0$
\end{opts4}
}

% 189 = 例5.8
\vbox{%
\examp{189.}{$f(x)=\displaystyle\int_{0}^{x}\dfrac{t}{t^{2}+2t+2}\,\mathrm{d}t$ 在 $[0,1]$ 上的最大值为 $\underline{\hspace{3em}}$.}
}

% 190 = 例15.5
\vbox{%
\examp{190.}{求定解问题 $\begin{cases} y''+2x(y')^{2}=0,\\ y(0)=1,\\ y'(0)=0\end{cases}$ 的解.}
}

% 191 = 例1.9
\vbox{%
\examp{191.}{$\lim\limits_{x\to\infty}\left[x\ln\left(\mathrm{e}+\dfrac{1}{x-1}\right)-x\right]=\underline{\hspace{5em}}$.}
}

% 192 = 例10.11
\vbox{%
\examp{192.}{计算下列曲线的弧长：\\
(1) $y=\ln\cos x$（$0\leqslant x\leqslant \dfrac{\pi}{6}$）；\\
(2) $\ln y+2x-\dfrac{1}{2}y^{2}=0$（$1\leqslant y\leqslant \mathrm{e}$）.}
}

% 193 = 例4.6
\vbox{%
\examp{193.}{设 $f(x)=\dfrac{1}{2^{x}+1}$，$x\in\mathbf{R}$，则 $f^{(4)}(0)=\underline{\hspace{3em}}$.}
}

% 194 = 例2.8
\vbox{%
\examp{194.}{设数列 $\{a_n\}$ 满足 $\begin{cases}a_1>2,\\ a_{n+1}=a_n-\dfrac{a_n-3}{a_n-2},\end{cases}$ $n=1,2,3,\cdots$，则（\quad）.}
\begin{opts2}
(A) $\{a_n\}$ 收敛于大于 $3$ 的数 & (B) $\{a_n\}$ 收敛于 $3$ \\
(C) $\{a_n\}$ 发散 & (D) $\{a_n\}$ 的敛散性与 $a_1$ 有关
\end{opts2}
}

% 195 = 例6.15
\vbox{%
\examp{195.}{设 $f(x)=\displaystyle\int_1^{x}\sin t^2\,\mathrm{d}t$，$g(x)=\displaystyle\int_0^{x}f(t)\,\mathrm{d}t$. \\[0.6ex]
(1) 计算 $g(1)$；\\[0.6ex]
(2) 证明：存在 $\xi\in(0,1)$，使得 $\displaystyle\int_1^{\xi}\sin t^2\,\mathrm{d}t=\xi(\cos 1-1)$.}
}

% 196 = 例6.7
\vbox{%
\examp{196.}{设函数 $f(x)$ 在 $[0,1]$ 上二阶可导，$f(0)=0$，且 $f(x)$ 在 $(0,1)$ 内取得最大值 $2$，在 $(0,1)$ 内取得最小值，证明：\\
(1) 存在 $\xi\in(0,1)$，使得 $f'(\xi)>2$；\\
(2) 存在 $\eta\in(0,1)$，使得 $f''(\eta)<-4$.}
}

% 197 = 例9.15
\vbox{%
\examp{197.}{设 $\displaystyle\int_1^{+\infty}\dfrac{a}{x(2x+a)}\,\mathrm{d}x=\ln 2$，则 $a=\underline{\hspace{3em}}$.}
}

% 198 = 例6.19
\vbox{%
\examp{198.}{设函数 $f(x)$ 在区间 $(a,b)$ 内可导. 证明：导函数 $f'(x)$ 在 $(a,b)$ 内严格单调增加的充分必要条件是对 $(a,b)$ 内任意的 $x_1$，$x_2$，$x_3$，当 $x_1<x_2<x_3$ 时，$\dfrac{f(x_2)-f(x_1)}{x_2-x_1}<\dfrac{f(x_3)-f(x_2)}{x_3-x_2}$.}
}

% 199 = 例4.7
\vbox{%
\examp{199.}{已知可导函数 $y=y(x)$ 满足 $a\mathrm{e}^{x}+y^{2}+y-\ln(1+x)\cos y+b=0$，且 $y(0)=0$，$y'(0)=0$. \\
(1) 求 $a,b$ 的值；\\
(2) 判断 $x=0$ 是否为 $y(x)$ 的极值点.}
}

% 200 = 例4.2
\vbox{%
\examp{200.}{(1) 设 $y=\dfrac{1}{x(1-x)}$，求 $\dfrac{\mathrm{d}^{n}y}{\mathrm{d}x^{n}}$；\\
(2) 设 $z=\dfrac{y^{2}}{x(1-x)}$，求 $\dfrac{\partial^{n}z}{\partial x^{n}}$.}
}

% 201 = 例11.16
\vbox{%
\examp{201.}{已知 $f'(\ln x)=\begin{cases}1, & x\in(0,1),\\ x, & x\in(1,+\infty),\end{cases}$ $f(0)=0$，则 $f(1)=\underline{\hspace{3em}}$.}
}

% 202 = 例7.2
\vbox{%
\examp{202.}{已知高温物体置于低温介质中，任一时刻该物体温度对时间的变化率与该时刻物体和介质的温差成正比. 现将一初始温度为 $120\,^\circ\mathrm{C}$ 的物体在 $20\,^\circ\mathrm{C}$ 恒温介质中冷却，$30\,\mathrm{min}$ 后该物体温度降至 $30\,^\circ\mathrm{C}$，若要将该物体的温度继续降至 $21\,^\circ\mathrm{C}$，还需冷却多长时间？}
}

% 203 = 例6.13
\vbox{%
\examp{203.}{已知 $\displaystyle\int_{0}^{x}f(t)\,\mathrm{d}t=xf(\theta x)$，$x>0$，$f(x)=\mathrm{e}^{x}$，则 $\lim\limits_{x\to 0^{+}}\theta=\underline{\hspace{3em}}$.}
}

% 204 = 例13.4
\vbox{%
\examp{204.}{求 $\lim\limits_{\substack{x\to+\infty\\ y\to 1}}(x^2+y^2)\mathrm{e}^{-(x+y)}$.}
}

% 205 = 例1.10
\vbox{%
\examp{205.}{$\lim\limits_{x\to+\infty}\left[\dfrac{x^{1+x}}{(1+x)^{x}}-\dfrac{x}{\mathrm{e}}\right]=\underline{\hspace{5em}}$.}
}

\end{document}
